\subsection{Opción 3 --- SCAN}

\begin{frame}{3 · Portada --- ¿Los modelos ``entienden'' o solo memorizan?}
  \OptionHeader{Recomendada nº 3}{Fundamentos / lenguaje}{%
    Medimos si un modelo que ``conoce'' una instrucción suelta es capaz de recombinarla en frases nuevas --- el famoso fallo tipo ``can't jump''.
  }
\end{frame}

\begin{frame}{3 · Imagina esto (sin jerga)}
  \begin{itemize}
    \item A un niño le enseñas el verbo nuevo ``dax'' y entiende enseguida ``dax dos veces''.
    \item Muchos modelos de IA, en cambio, fallan en esa recombinación.
    \item No es un debate filosófico vacío: se puede \textbf{medir} con un examen diseñado para eso.
    \item SCAN es ese examen: órdenes simples (``salta'', ``anda'', ``dos veces'') convertidas en acciones.
  \end{itemize}
\end{frame}

\begin{frame}{3 · Analogía cotidiana}
  \begin{center}
  \begin{tikzpicture}
    \node[box, text width=2.6cm] (a) {Sabes\\``saltar''};
    \node[box, text width=2.6cm, right=0.45cm of a] (b) {Sabes\\``dos veces''};
    \node[boxdark, text width=3.0cm, right=0.45cm of b] (c) {¿Entiendes\\``salta 2 veces''?};
    \draw[arr] (a)--(b)--(c);
  \end{tikzpicture}
  \end{center}
  \vspace{0.45em}
  Si solo memorizaste ejemplos raros del examen, puedes aprobar el tipo A y suspender el tipo B.\\
  Eso es \textbf{generalización composicional}.
\end{frame}

\begin{frame}{3 · ¿Qué haría aquí el machine learning?}
  \begin{itemize}
    \item Entrenar modelos clásicos y redes en un trozo de SCAN.
    \item Probarlos en el split donde ``jump'' casi no aparece compuesto.
    \item Comparar: ¿quién falla y en qué tipo de órdenes?
    \item Opcional: repetir la idea en un dataset más lingüístico (COGS) para responder ``¿SCAN ya está resuelto?''.
  \end{itemize}
\end{frame}

\begin{frame}{3 · Glosario rápido}
  \begin{description}[font=\bfseries\scriptsize,style=nextline,leftmargin=1.5cm]
    \item[Composición] Combinar piezas conocidas en algo nuevo.
    \item[OOD] Fuera de lo visto en entrenamiento (el examen ``difícil'').
    \item[SCAN] Banco de pruebas de órdenes $\rightarrow$ acciones (Lake \& Baroni).
    \item[addprim\_jump] Split concreto donde ``jump'' solo se ve aislado en train.
  \end{description}
\end{frame}

\begin{frame}{3 · Pregunta de investigación}
  \begin{block}{Pregunta}
    Cuando el modelo ve ``jump'' solo en aislamiento, ¿qué familias de métodos fallan al usarlo dentro de órdenes nuevas, y qué señales del input predicen ese fallo?
  \end{block}
\end{frame}

\begin{frame}{3 · Cómo lo haríamos, paso a paso}
  \PipelineFive{Splits\\oficiales}{Modelos\\clásicos}{Redes\\secuencia}{Análisis\\de errores}{COGS\\opcional}
\end{frame}

\begin{frame}{3 · Datos}
  \begin{itemize}
    \item SCAN: pequeño, limpio, se corre en un portátil.
    \item Paper clásico (ICML 2018) = bibliografía fuerte.
    \item Defensa ante ``¿es realista?'': es el banco de pruebas estándar de composición, como un test de conducción controlado.
  \end{itemize}
\end{frame}

\begin{frame}{3 · Métodos del curso y más allá}
  \MethodsSplit{N-gramas / TF-IDF · Logística / SVM · Árboles · Validación OOD}{LSTM/GRU · Transformer pequeño · Análisis de errores · Modelo de lenguaje abierto solo como contraste}
\end{frame}

\begin{frame}{3 · Por qué encaja con el curso}
  \begin{columns}[T]
    \begin{column}{0.55\textwidth}
      \begin{itemize}
        \item Muy original frente a Kaggle típico.
        \item La comparación de métodos es el núcleo.
        \item Oral memorable si se cuenta bien.
      \end{itemize}
    \end{column}
    \begin{column}{0.42\textwidth}
      \ScoreBar{Ambición}{9}{ieblue}
      \ScoreBar{Factibilidad}{9}{good}
      \ScoreBar{Diferenciación}{10}{accent}
      \ScoreBar{Portfolio VC}{5}{bad}
    \end{column}
  \end{columns}
\end{frame}

\begin{frame}{3 · Por qué encaja con nosotros}
  \FitDual{Quiere profundidad tipo investigación y temas de abducción / límites de los LLM.}{Le gusta NLP y comparar modelos clásicos frente a redes con análisis claro.}{10}{8}
\end{frame}

\begin{frame}{3 · Pros y contras}
  \ProCon{Muy diferenciador\\Muy factible\\Oral fuerte\\Encaja con los dos}{Datos sintéticos\\Crítica ``ya está estudiado''\\Menos portfolio espacial/VC\\Conviene añadir COGS o diagnóstico}
\end{frame}

\begin{frame}{3 · Impacto}
  Aclarar cuándo un modelo parece inteligente pero no recombina reglas --- útil para diseñar evaluaciones honestas.
\end{frame}

\begin{frame}{3 · Veredicto}
  \Verdict{Muy recomendada (prioridad alta)}{Ideal si la identidad es ``fundamentos / composición''. Alternativa cercana: SCAR (nº 8).}
\end{frame}
